\section{System model}
\label{sec:model}

\begin{inbrief}
The platform is a central computer and zone controllers with typed
cores, a shared memory per ECU, and a TSN network (\cref{sec:platform,sec:network-model}); \cref{sec:other-architectures} extends it to
domain, hybrid and distributed architectures and to any network topology. The application is a set of multi-rate task-chains
with bindings to ECUs and safety constraints (\cref{sec:application}). Cores
communicate at LET instants, with implicit communication or strict LET (\cref{sec:execution}). \Cref{sec:metrics}
defines how a mapped task-chain is cut into segments and inter-ECU
communications, and which timing metrics are computed.
\Cref{tab:notation} collects the notation.
\end{inbrief}

Times are integers in microseconds and utilizations are fractions of a
core. The implementation stores utilizations as integers in parts per
million and rounds them up, so every rounding is conservative.

\begin{table}[t]
\centering\footnotesize
\caption{Main notation.}
\label{tab:notation}
\setlength{\tabcolsep}{3pt}
\rowcolors{2}{brandlight}{white}
\begin{tabularx}{\columnwidth}{@{}lX@{}}
\toprule
\rowcolor{white}
symbol & meaning \\
\midrule
$\Ctrl=\{k_0,\dots,k_Z\}$ & ECUs; on the zonal platform, $k_0$ the central computer, $k_1..k_Z$ the zone controllers \\
$\Types$, $\Types(k)$ & core types; types available on ECU $k$ \\
$s_\theta$, $s_{k,\theta}$ & speed of type $\theta$ (on ECU $k$), \% of the reference core \\
$\Cores$, $\kappa(p)$, $\theta(p)$ & cores; ECU and type of core $p$ \\
$\varphi$ & utilization cap of a core ($0<\varphi\le1$) \\
$\Chains$, $\xi$ & application: set of task-chains; one task-chain \\
$\tau$, $T_\tau$, $C^{\mathrm{ref}}_\tau$ & task, its period, its WCET on the reference core \\
$U^{\mathrm{ref}}_\tau$ & utilization on the reference core, $C^{\mathrm{ref}}_\tau/T_\tau$ \\
$f_{\tau,\theta}$ & WCET factor of $\tau$ on type $\theta$ (\%) \\
$C^{\mathrm{iso}}_\tau(p)$, $C_\tau(p)$, $U_\tau(p)$ & WCET of $\tau$ on core $p$ alone, with contention; utilization \\
$a_\tau$, $\ell$, $n_k$ & memory requests per job; time per request; competing cores of ECU $k$ \\
$A_\tau$, $\Theta_\tau$ & ECUs and core types $\tau$ may run on \\
$\mu_\tau$, $M_k$ & memory of $\tau$; memory of ECU $k$ \\
$D_\xi$ & end-to-end deadline of $\xi$ (on its maximum reaction time) \\
$\rho(\xi)$ & replica group of $\xi$ \\
$\Frags$, $\pi$ & sub-chains; mapping of sub-chains to cores \\
$T^{\LET}_p$, $C^{\LET}_p$ & period and WCET of the LET task of core $p$ \\
$R_\tau$ & worst-case response time (WCRT) of $\tau$ \\
$\sigma$ & utilization limit of a sub-chain (fraction of a core) \\
$G=(V,E)$, $r_e$ & network graph; rate of link $e$ (bit/\si{\micro s}) \\
$\delta$ & latency of a switch \\
$X$, $\vartheta_x$, $l_x$ & traffic classes; period threshold of class $x$; its largest frame at a port \\
$b_h$, $J_h$, $T_h$ & bits, release jitter and period of flow $h$ \\
$\sigma_h$, $\rho_h$ & burst and rate of flow $h$ (unrelated to $\sigma$ and $\rho(\xi)$) \\
$d_{e,x}$, $d_h$ & delay bound of class $x$ at port $e$; of flow $h$ \\
$w_h$ & copy-out delay of communication $h$ (\cref{eq:copy-out}) \\
$\Lambda_\xi$, $\Gamma_\xi$, $\Delta_\xi$ & end-to-end response time (holistic sum), end-to-end delay, maximum reaction time of $\xi$ \\
$M(\pi)$, $N(\pi)$ & cores used; task-chains on time \\
$\lambda(\pi)$, $\lambda^\star$ & critical scaling factor of a mapping; its target \\
\bottomrule
\end{tabularx}
\end{table}

\subsection{Platform}
\label{sec:platform}

\paragraph{ECUs and cores.} The platform has a set of ECUs $\Ctrl=\{k_0,
k_1,\dots,k_Z\}$ (\cref{fig:architecture}). $k_0$ is the central computer
and $k_1,\dots,k_Z$ are the zone controllers; the default is $Z=5$. An ECU
holds cores of the types $\Types(k)$ it supports. The central computer
offers \emph{big} application cores and, optionally, slower \emph{little}
cores. A zone controller offers microcontroller cores (\emph{MCU}) and,
optionally, \emph{lock-step} cores (LS), a common safety mechanism for
the functions of high integrity levels under ISO~26262 \cite{iso26262}. A type $\theta$ has a speed $s_\theta$, in percent of
a reference core. By default $s_{\mathrm{big}}=100$, $s_{\mathrm{little}}
=40$ and $s_{\mathrm{MCU}}=s_{\mathrm{LS}}=50$.

The initial architecture $\Cores_0$ has a number of cores of each type on
each ECU: by default 8 big cores on the central computer and 4 MCU cores on
each zone controller, 28 cores in all. When the application needs more,
cores are added. The number of cores \emph{used}, that is, running at least
one application task, is the main sizing result, $M$.

\paragraph{Utilization cap and memory.} The total utilization of the
tasks of a core, its LET task included, may not exceed the utilization cap
$\varphi$ (default 0.8). The remaining $1-\varphi$ is the design reserve
for estimation errors and future functions. ECU $k$ has a memory $M_k$ for
the code and data of its tasks (unlimited by default).

\begin{figure}[t]
\centering
\resizebox{\columnwidth}{!}{%
\begin{tikzpicture}[x=1cm,y=1cm,
	core/.style={draw=cmuted, fill=white, minimum size=3.2mm,
		inner sep=0pt},
	ctrl/.style={draw=cink, rounded corners=2pt, inner sep=2pt},
	dev/.style={draw=cmuted, circle, inner sep=0.6pt, font=\tiny\sffamily,
		fill=cyellow!25}]
\foreach \i in {0,...,3} \node[core] (p\i) at (\i*0.4,0) {};
\foreach \i in {4,...,7} \node[core] (p\i) at ({(\i-4)*0.4},-0.4) {};
\node[ctrl, fill=cpcu, fit=(p0)(p7),
	label={[font=\scriptsize\sffamily]above:central computer: big cores (+ little)}] (pcu) {};
\foreach \i in {0,...,7} \node[core] at (p\i) {};
\node[draw=cink, fill=cnet, minimum width=1.3cm, minimum height=0.45cm,
	font=\scriptsize\sffamily] (sw) at (0.6,-1.55) {TSN switch};
\draw[very thick, color=corange] (pcu) -- (sw);
\foreach \z/\x/\y in {1/-2.9/-1.9, 2/-1.75/-3.0, 3/0.6/-3.35, 4/2.95/-3.0, 5/4.1/-1.9} {
	\node[core] (z\z a) at (\x-0.2,\y+0.1) {};
	\node[core] (z\z b) at (\x+0.2,\y+0.1) {};
	\node[core] (z\z c) at (\x-0.2,\y-0.3) {};
	\node[core] (z\z d) at (\x+0.2,\y-0.3) {};
	\node[ctrl, fill=czcu, fit=(z\z a)(z\z d),
		label={[font=\scriptsize\sffamily]below:ZCU\z}] (zcu\z) {};
	\foreach \c in {a,b,c,d} \node[core] at (z\z\c) {};
	\draw[thick, color=corange] (zcu\z) -- (sw);
}
\node[dev] at ($(zcu1)+(-0.75,0.35)$) (s1) {S};
\node[dev] at ($(zcu1)+(-0.75,-0.3)$) (a1) {A};
\draw[cmuted] (s1) -- (zcu1); \draw[cmuted] (a1) -- (zcu1);
\node[dev] at ($(zcu5)+(0.75,0.35)$) (s5) {S};
\node[dev] at ($(zcu5)+(0.75,-0.3)$) (a5) {A};
\draw[cmuted] (s5) -- (zcu5); \draw[cmuted] (a5) -- (zcu5);
\node[dev] at ($(zcu3)+(-0.75,-0.25)$) (s3) {S};
\draw[cmuted] (s3) -- (zcu3);
\node[note, anchor=west, text width=2.3cm, align=left] at (2.35,-0.2)
	{thick: backbone rate\\ S/A: sensors and actuators wired to a zone};
\end{tikzpicture}}
\caption{The zonal platform (star topology, default sizes): a central
computer, five zone controllers (ZCU) of four cores, one TSN switch. Cores
are added where the application needs them.}
\label{fig:architecture}
\end{figure}

\paragraph{Execution times.} Every task has a reference WCET
$C^{\mathrm{ref}}_\tau$, measured on the reference core, and a factor
$f_{\tau,\theta}$ (in percent, default 100) for each core type. The factor
models tasks that gain or lose more than the speed ratio on a type, for
instance through caches or vector units. On a core $p$ of type $\theta=
\theta(p)$, the WCET in isolation is
\begin{equation}
C^{\mathrm{iso}}_\tau(p)=\max\!\left(1,\ceil{\frac{C^{\mathrm{ref}}_\tau\,
f_{\tau,\theta}}{s_\theta}}\right).
\label{eq:wcet-iso}
\end{equation}

\paragraph{Shared-memory contention.}
\label{sec:contention}
The cores of one ECU share its memory. A job of task $\tau$ issues at most
$a_\tau$ requests to it. The memory serves one request in $\ell$
nanoseconds, and serves the cores in round robin. Each core has at most one
request pending at a time (in-order cores, blocking memory accesses). A
request therefore waits at most for one request of every other core that
competes for the memory \cite{paolieri2009hardware}. If $n_k$ cores of ECU $k$ compete, the
WCET of $\tau$ on a core $p$ of $k$ is
\begin{equation}
C_\tau(p)=C^{\mathrm{iso}}_\tau(p)+\ceil{\frac{a_\tau\,\ell\,(n_k-1)}{1000}},
\qquad U_\tau(p)=\frac{C_\tau(p)}{T_\tau}.
\label{eq:wcet}
\end{equation}
The waiting time does not depend on the speed of the core, which is why it
is added after the scaling of \cref{eq:wcet-iso}. The number of competing
cores depends on the mapping: \cref{sec:contention-method} explains how
the method fixes it. When no request count is given, a task gets $a_\tau=
\lceil\eta\,C^{\mathrm{ref}}_\tau/1000\rceil$ requests, a density of
$\eta$ requests per millisecond of reference WCET (default 1000). With
$\ell=0$, the default, there is no contention and $C_\tau(p)=
C^{\mathrm{iso}}_\tau(p)$.

\subsection{Application}
\label{sec:application}

\begin{definition}[Task]
A task $\tau$ is periodic with period $T_\tau$ and an implicit deadline
equal to it. Its attributes are the reference WCET $C^{\mathrm{ref}}_\tau$
and the type factors $f_{\tau,\theta}$; the number of memory requests
$a_\tau$; the size $m_\tau$ (bytes) of the output it produces at every
job; its memory $\mu_\tau$; the set $A_\tau\subseteq\Ctrl$ of ECUs and the
set $\Theta_\tau\subseteq\Types$ of core types it may run on; and a
co-location flag $g_\tau$. With $g_\tau$ set, the task must run on the
same core as its successor in the chain.
\end{definition}

A task with $|A_\tau|=1$ is \emph{pinned} to its ECU: sensor and actuator
drivers run on the zone controller their device is wired to. The other
tasks are \emph{free}. A free task may still be restricted, for instance
to the zone controllers only ($A_\tau=\{k_1,\dots,k_Z\}$), or kept out of
the zones of its replica (below).

\begin{definition}[Task-chain]
A task-chain $\xi=(\tau_1,\dots,\tau_n)$ is a sequence of tasks in which
the output of $\tau_i$ is the input of $\tau_{i+1}$. A maximal run of
consecutive tasks with the same period is an \emph{activation pattern}
\cite{kramer2015waters}. If two patterns of a task-chain have the same
period, their tasks on a common core are also kept in chain order, which
only restricts the priorities. The task-chain has an end-to-end deadline $D_\xi$
on its maximum reaction time (\cref{sec:metrics}) and, possibly, a replica
group $\rho(\xi)$.
\end{definition}

The links of a task-chain come in two kinds (\cref{fig:chain}).
\begin{itemize}
\item \emph{Inside a pattern}, tasks are released together. On one core
they execute in chain order: there is a \emph{precedence} constraint, so
a job uses the data of the job of its predecessor released with it. When
the partitioning cuts a pattern, for instance at a sensor driver bound to
a zone, the two parts run on different cores, and the link becomes a
communication: the consumer then reads the data of an earlier job, with up
to one consumer period of sampling, which \cref{sec:e2e} bounds.
\item \emph{Between two patterns}, data goes through a \emph{register}:
the consumer reads the latest value at its own release. There is no
precedence, but there is a sampling delay of up to one consumer period.
\end{itemize}

The task-chains of one replica group are redundant implementations of one
safety function, with redundant sensors and actuators. No two of them may
run on the same ECU or, under a weaker policy, on the same core
\cite{bhat2017practical}.

\begin{figure}[t]
\centering
\begin{tikzpicture}[node distance=1.6mm]
\node[task, fill=czcu] (t1) {$\tau_1$};
\node[task, right=of t1] (t2) {$\tau_2$};
\node[task, right=of t2] (t3) {$\tau_3$};
\node[task, right=5mm of t3] (t4) {$\tau_4$};
\node[task, right=of t4] (t5) {$\tau_5$};
\node[task, right=5mm of t5] (t6) {$\tau_6$};
\node[task, right=of t6, fill=czcu] (t7) {$\tau_7$};
\foreach \a/\b in {t1/t2,t2/t3,t4/t5,t6/t7}
	\draw[arrow, color=cink] (\a) -- (\b);
\draw[arrow, dashed, color=corange] (t3) -- (t4);
\draw[arrow, dashed, color=corange] (t5) -- (t6);
\draw[decorate, decoration={brace, amplitude=3pt, mirror}]
	([yshift=-1.5mm]t1.south west) -- node[note, below=2pt] {$T=10$\,ms} ([yshift=-1.5mm]t3.south east);
\draw[decorate, decoration={brace, amplitude=3pt, mirror}]
	([yshift=-1.5mm]t4.south west) -- node[note, below=2pt] {$T=20$\,ms} ([yshift=-1.5mm]t5.south east);
\draw[decorate, decoration={brace, amplitude=3pt, mirror}]
	([yshift=-1.5mm]t6.south west) -- node[note, below=2pt] {$T=5$\,ms} ([yshift=-1.5mm]t7.south east);
\node[note, above=1.5mm of t1] {sensor, ZCU2};
\node[note, above=1.5mm of t7] {actuator, ZCU4};
\node[note, above=1.5mm of t4, xshift=3mm] {free tasks};
\end{tikzpicture}
\caption{A task-chain of three activation patterns. Solid arrows:
precedence inside a pattern; dashed arrows: registers between patterns.
Shaded tasks are pinned to a zone controller.}
\label{fig:chain}
\end{figure}

\subsection{Scheduling and communication on the cores}
\label{sec:execution}

\paragraph{Scheduling.} Every core schedules its tasks with preemptive
fixed priorities. Inside a pattern, the tasks mapped to the same core get
priorities that decrease along the chain, which enforces the precedence on
that core.

\paragraph{Communication between cores.} Tasks on different cores never
share data directly. Every core $p$ that runs at least one application
task runs one \emph{LET task}, a copy task with the highest priority on
$p$, released at the \emph{LET instants} of $p$. It copies data between
the local memory of $p$ and the global memory, or the network interface,
and
\begin{gather}
T^{\LET}_p=\gcd\{T_\tau:\tau\text{ on }p\},\nonumber\\
C^{\LET}_p=\ceil{\frac{100\,c_{\LET}}{s_{\theta(p)}}}
	+\ceil{\frac{a_{\LET}\,\ell\,(n_k-1)}{1000}},
\label{eq:let}
\end{gather}
so that every release of every task of $p$ is a LET instant. A consumer
on another core reads, at its release, the last value copied in by then.
Two schemes decide when an output is copied out.
\begin{itemize}
\item \emph{Implicit communication} (the default): an output is copied at
the first LET instant after its producer finishes. It is the implicit
communication of AUTOSAR \cite{hamann2017communication}, with the copies
grouped at the LET instants. The delays are short, but the data flow
depends on response times.
\item \emph{Strict LET} \cite{kirsch2012let,biondi2018achieving}: an output
is published at the end of the period of its producer, which is a LET
instant of its core, whatever the response time. The data flow no longer
depends on response times, at the price of longer delays.
\end{itemize}
\Cref{sec:e2e} bounds both, and \cref{sec:rq3} compares them. Here
$c_{\LET}$ is the copy time on the reference core (default
\SI{50}{\micro s}, of the order of the LET communication overheads that
Biondi and Di Natale measured on a multicore \cite{biondi2018achieving}),
and $a_{\LET}=\lceil\eta\,c_{\LET}/1000\rceil$ are its memory requests.
The copy time is a constant per LET instant, whatever the amount of data
copied (\cref{sec:discussion}).
The LET load $C^{\LET}_p/T^{\LET}_p$ belongs to the core. On a reference
core, a 1\,ms task anywhere on it costs 5\,\% of LET load, while periods that
are all multiples of 10\,ms cost 0.5\,\%. The period mix of a core
therefore matters, which is one reason why mapping is more than packing
utilizations.

The \emph{load} of core $p$ under a mapping is
\begin{equation}
\mathcal{L}(p)=\sum_{\tau\text{ on }p} U_\tau(p) + \frac{C^{\LET}_p}{T^{\LET}_p}
\;\le\;\varphi .
\label{eq:load}
\end{equation}

\subsection{Network}
\label{sec:network-model}

\paragraph{Topology and routing.} The network is a directed graph $G=(V,E)$
whose nodes are the ECUs and the switches. Every full-duplex cable gives
two directed links, and each link $e$ has an output port at its origin and
a rate $r_e$. Three topologies are considered (\cref{fig:topologies}):
\begin{itemize}
\item \emph{star}: one switch, one link per ECU;
\item \emph{tree}: one switch per ECU, and the zone switches linked to the
switch of the central computer;
\item \emph{ring}: one switch per ECU, the switches in a ring.
\end{itemize}
ECU links run at the link rate $r$ (default 1\,Gbit/s). Links between
switches, and the link of the central computer, which carries most of the
traffic, run at the backbone rate $r_b\ge r$ (by default $r_b=r$). A
message follows a shortest path. ECUs do not forward, and ties are broken
by a fixed link order (clockwise on the ring). Every switch adds a
constant latency $\delta$ (default \SI{2}{\micro s}) after
store-and-forward.

\begin{figure}[t]
\centering
\begin{tikzpicture}[
	c/.style={draw=cink, rounded corners=1pt, minimum size=3.3mm,
		inner sep=0pt, font=\tiny\sffamily, fill=czcu},
	p/.style={c, fill=cpcu},
	s/.style={draw=cink, circle, minimum size=2.6mm, inner sep=0pt,
		fill=cnet},
	l/.style={color=corange, semithick},
	b/.style={color=corange, very thick}]
\begin{scope}
\node[s] (S) at (0,0) {};
\node[p] (P) at (0,0.7) {C};
\foreach \i/\a in {1/200,2/245,3/295,4/340,5/20}
	\node[c] (Z\i) at (\a:0.75) {\i};
\draw[b] (P) -- (S);
\foreach \i in {1,...,5} \draw[l] (Z\i) -- (S);
\node[note] at (0,-1.22) {star};
\end{scope}
\begin{scope}[xshift=2.65cm]
\node[s] (S0) at (0,0.25) {};
\node[p] (P) at (0,0.8) {C};
\draw[b] (P) -- (S0);
\foreach \i/\x in {1/-1.0,2/-0.5,3/0,4/0.5,5/1.0} {
	\node[s] (S\i) at (\x,-0.25) {};
	\node[c] (Z\i) at (\x,-0.68) {\i};
	\draw[l] (Z\i) -- (S\i);
	\draw[b] (S\i) -- (S0);
}
\node[note] at (0,-1.22) {tree};
\end{scope}
\begin{scope}[xshift=5.4cm]
\foreach \i/\a in {0/90,1/150,2/210,3/270,4/330,5/30}
	\node[s] (S\i) at (\a:0.45) {};
\foreach \i/\j in {0/1,1/2,2/3,3/4,4/5,5/0} \draw[b] (S\i) -- (S\j);
\node[p] (C0) at (90:0.88) {C};
\draw[b] (C0) -- (S0);
\foreach \i/\a in {1/150,2/210,3/270,4/330,5/30} {
	\node[c] (C\i) at (\a:0.88) {\i};
	\draw[l] (C\i) -- (S\i);
}
\node[note] at (0,-1.22) {ring};
\end{scope}
\end{tikzpicture}
\caption{Topologies. C: central computer; 1--5: zone controllers; circles:
switches. Thick links run at the backbone rate $r_b$, thin ones at $r$.}
\label{fig:topologies}
\end{figure}

\paragraph{Frames.} A message of $m$ bytes is cut into frames of at most
1500 bytes of payload. Each frame carries 42 bytes of overhead: header 14,
VLAN tag 4, checksum 4, preamble 8 and inter-frame gap 12. A payload is
padded to at least 42 bytes, the 64-byte minimum frame with a VLAN tag.
The bits on the wire are
\begin{equation}
b(m)=\sum_{i}8\,\bigl(\max(m_i,42)+42\bigr),
\label{eq:bits}
\end{equation}
over the payload chunks $m_i$ of the message.

\paragraph{Traffic classes.} LET traffic uses $X\in\{1,2,3\}$ classes of
IEEE 802.1Q \cite{ieee8021q2022}. A stream whose producer has period $T$
belongs to the first class $x$ with $T\le\vartheta_x$; class 0 has the
highest priority. The default thresholds are 20\,ms for $X=2$, and
5\,ms and 20\,ms for $X=3$. Other traffic is best effort. It is assumed to
be always present, with maximum-size frames of $L_{BE}=1542$ bytes (12\,336 bits) on the
wire.

\paragraph{Shaping.} Every port shapes the LET classes with one of three
mechanisms (\cref{sec:background}).
\begin{itemize}
\item \emph{CBS}: one credit-based shaper per class, with idle slopes
$I_{e,x}$ taken from a budget of shares of $r_e$ (default 75\,\% for one
class; 50 and 25\,\% for two; 50, 15 and 10\,\% for three).
\item \emph{TAS}: one exclusive window per class in a gate cycle $T_c$
(default \SI{500}{\micro s}), with guard bands (default windows 70\,\%;
40 and 30\,\%; 30, 20 and 20\,\%).
\item \emph{ATS}: strict priority between the classes, and at every switch
one interleaved regulator per input port and class. Each stream is
regulated to its token bucket at its source (\cref{sec:nc}).
\end{itemize}
With frame preemption (802.1Qbu), a best-effort frame blocks for at most
143 bytes \cite{thiele2016formal}. A link may also carry a design limit on
its LET load (default 100\,\%).

\paragraph{Frame replication.} On the ring, frame replication (FRER) may
be enabled for the task-chains of the replica groups, whose redundancy
would otherwise end at the network. Each of their inter-ECU messages is
then sent as two member streams (\cref{fig:frer}). The switch of the source
ECU replicates every frame onto both directions of the ring. The switch of
the destination ECU keeps the first copy of each sequence number and
discards the other. The two routes share only the links between the ECUs
and their switches, so a single failure of a link between switches, or of
a switch other than these two, leaves one member stream.

\begin{figure}[t]
\centering
\begin{tikzpicture}[
	s/.style={draw=cink, circle, minimum size=3.2mm, inner sep=0pt,
		fill=cnet, font=\tiny\sffamily},
	c/.style={draw=cink, rounded corners=1pt, minimum size=3.6mm,
		inner sep=0pt, font=\tiny\sffamily, fill=czcu}]
\foreach \i/\a in {0/90,1/150,2/210,3/270,4/330,5/30}
	\node[s] (S\i) at (\a:0.95) {};
\foreach \i/\j in {0/1,1/2,2/3,3/4,4/5,5/0} \draw[cmuted] (S\i) -- (S\j);
\node[c] (Z1) at (150:1.55) {1};
\node[c] (Z2) at (210:1.55) {2};
\draw[very thick, color=cink] (Z1) -- (S1);
\draw[very thick, color=cink] (S2) -- (Z2);
\draw[-{Stealth[length=1.6mm]}, very thick, color=cblue] (S1) -- (S2);
\draw[-{Stealth[length=1.6mm]}, very thick, color=corange, dashed]
	(S1) -- (S0) -- (S5) -- (S4) -- (S3) -- (S2);
\node[note, anchor=west, align=left] at (1.35,0.55)
	{\textcolor{cblue}{member stream 1: short route}\\
	 \textcolor{corange}{member stream 2: other way}};
\node[note, anchor=west, align=left] at (1.35,-0.35)
	{replicated in the switch of ZCU1,\\ eliminated in the switch of ZCU2};
\end{tikzpicture}
\caption{Frame replication on the ring for a message from ZCU1 to ZCU2.
The two member streams share only the two links between the ECUs and their
switches.}
\label{fig:frer}
\end{figure}

\subsection{Other E/E architectures}
\label{sec:other-architectures}

The zonal platform above is the running example, but nothing in the
model needs it. In general, a platform is a list of ECUs, and each ECU
$k$ has a \emph{role}: central computer, zone controller, domain
controller, or ECU of a distributed architecture. It also has initial
cores of given types, a speed $s_{k,\theta}$ for each of its types
(the speed of a type may differ from one ECU to another), a memory $M_k$,
the type of the cores added to it (another of its types is added for a
sub-chain that may not run on that one), and a largest number of
cores $M^{\max}_k$ (unbounded by default). The role only matters for the
sub-chains that are bound to no ECU (\emph{free} sub-chains): they go to
their \emph{hosts} first, and the cores added for them are added to a
host first. The hosts are the central computers. On a platform without a
central computer, they are the ECUs whose fastest initial cores are the
fastest of the platform, where a sub-chain costs the least utilization
(\cref{sec:rq1} measures what this choice is worth). A free sub-chain
tries the cores of the hosts first, then the initial cores of the other
ECUs it is allowed on; it uses a core added to a non-host ECU only if no
host admits it (the communication-aware rule \CADU\ of \cref{sec:method}
weighs the hosts against the hops instead). A new
core for it goes to the first ECU that admits it, the hosts first and then
the others, each in the order of the platform, provided the ECU has fewer
than $M^{\max}_k$ cores, enough memory and no replica of its task-chain.
Everything else (allowed ECUs, core types,
memory, replica separation, schedulability, LET, the network analysis)
is defined per ECU and per core, and does not depend on the role.

The network may also be any set of switches and full-duplex links
(\emph{custom} topology), each ECU attached to one switch, with a rate per
link. Routes are shortest paths. With FRER, the second member stream of a
replicated flow takes the shortest route that shares no link and no
intermediate switch with the first; a flow without such a route is not
replicated. On the ring this is the other way around the ring, as above.

\Cref{tab:templates} lists the four platforms of the evaluation
(\cref{sec:rq7}): the zonal platform, a domain architecture of five
domain controllers, a hybrid one (a central computer, an ADAS domain
controller and four zone controllers), and a distributed one of twelve
small ECUs.

\begin{table}[t]
\centering\footnotesize
\caption{The four platforms of the evaluation: initial cores per ECU
(B: big, L: little at 40\,\%, M: MCU at 50\,\%, LS: lock-step at
50\,\%), star topology. Cores are added of the first type of each ECU.}
\label{tab:templates}
\setlength{\tabcolsep}{3pt}
\rowcolors{2}{brandlight}{white}
\begin{tabularx}{\columnwidth}{@{}lX@{}}
\toprule
\rowcolor{white}
platform & ECUs and initial cores \\
\midrule
zonal & central computer 8\,B; five zone controllers of 4\,M; 1\,Gbit/s \\
domain & ADAS 4\,B; infotainment 2\,B + 4\,L; powertrain 3\,M + 1\,LS;
chassis 2\,M + 2\,LS; body 2\,M; 1\,Gbit/s \\
hybrid & central computer 8\,B; ADAS domain controller 4\,B; four zone
controllers of 4\,M; 1\,Gbit/s \\
distributed & ADAS 2\,B (at most 4); ESC and EPS 1\,LS each; nine ECUs
of 1\,M (engine, transmission, body, climate, cluster, four doors); at
most 2 cores per ECU except ADAS; 100\,Mbit/s \\
\bottomrule
\end{tabularx}
\end{table}

\subsection{Mapping, segments and timing metrics}
\label{sec:metrics}

A task-chain is mapped as a sequence of \emph{sub-chains}: runs of
consecutive tasks placed together on one core. Following the chain,
consecutive sub-chains on the same core form a \emph{segment}. Between two
segments, data is sent either inside an ECU (shared memory) or between two
ECUs, over the network. Every \emph{inter-ECU communication} $h$ is a
network \emph{flow}. It carries one message of $b_h=b(m)$ bits per period
$T_h$ of its producer (the last task of the source segment), released at a
LET instant of the source core: after the producer has finished
(implicit communication) or at the end of its period (strict LET).

Three end-to-end quantities are used for a task-chain $\xi$
\cite{feiertag2008compositional,becker2017end,durr2019end}.
\begin{itemize}
\item The \emph{maximum reaction time} $\Delta_\xi$ is the longest time
from an external event to the end of the first job of $\tau_n$ that
accounts for it, every sampling delay included. It equals the maximum
data age \cite{gunzel2023equivalence}, and it is the metric compared with
the deadline.
\item The \emph{end-to-end delay} $\Gamma_\xi$ is the longest time from the
release of a job of $\tau_1$ to the end of the first job of $\tau_n$ that
uses its data, every copy and sampling delay included. It is the reaction
time without the wait for the first sample.
\item The \emph{end-to-end response time} $\Lambda_\xi$ is the holistic
sum \cite{tindell1994holistic}: the response times of the runs of the
chain (below) plus the network delays, without any copy or sampling
delay. It is not a time of the system but a figure of merit, which shows
how much of $\Delta_\xi$ is computation and network.
\end{itemize}
A task-chain \emph{meets its deadline} when all its tasks meet theirs,
every port its flows cross has a delay bound, and $\Delta_\xi\le D_\xi$.
\Cref{sec:e2e} bounds $\Gamma_\xi$ and $\Delta_\xi$.
